\section{Experiments and validation}
\label{sec:experiments}

\paragraph{Methodology.} Because parcel spawns are random, every
(scenario,~strategy) pair is run five times and compared by mean score. The
map is selected server-side per run; a campaign runner starts and stops the
server for each map, runs every strategy with a fresh in-game identity (so
scores start at 0 and never bleed across runs), and an aggregator produces the
comparison tables. The Challenge~1 campaign covered all eight \texttt{26c1\_*}
maps, four strategies, five 120\,s runs each (160 runs).

\paragraph{Baseline result.} \texttt{greedy-nearest} is the strongest strategy:
it wins outright on six maps and is statistically tied on the other two
(\texttt{26c1\_3}, \texttt{26c1\_5}, where all four strategies remain within
roughly 20--25 points, small relative to the run-to-run variance of the $n{=}5$
campaign), so it is never \emph{significantly} beaten. The best non-greedy
strategy on every map is \texttt{delivery-threshold} (\texttt{d.thr.}\ in the
table).

\begin{center}\small
\begin{tabular}{@{}lrrl@{}}
\toprule
Map & greedy & best other & outcome \\
\midrule
26c1\_1 & 868 & 742 (d.thr.) & greedy \\
26c1\_2 & 866 & 691 (d.thr.) & greedy \\
26c1\_3 & 377 & 397 (d.thr.) & tie (within noise) \\
26c1\_4 & 828 & 657 (d.thr.) & greedy \\
26c1\_5 & 319 & 337 (d.thr.) & tie (within noise) \\
26c1\_6 & 917 & 622 (d.thr.) & greedy \\
26c1\_7 & 882 & 290 (d.thr.) & greedy \\
26c1\_8 & 1035 & 604 (d.thr.) & greedy \\
\bottomrule
\end{tabular}
\end{center}

\paragraph{Diagnosis.} A log-level analysis explains \emph{why}. On the large
maps \texttt{greedy} delivers about 7.6 parcels per drop-off because it delivers
only when no parcel is reachable, hoarding on these parcel-rich maps, whereas
the value-aware strategies deliver in $\sim$2-parcel batches and so collect far
fewer parcels (on 26c1\_7, 33.8 vs.\ 10.6 picked up). Crucially there are
\emph{no real plan failures} for either --- the logged failures are all
intention-revision cancellations --- so this is a strategy/utility effect, not
a pathfinding bug. This \textbf{falsified} our prior hypothesis that the
slow-decay, high-reward maps (26c1\_7/8) would favour batching: there
\texttt{delivery-threshold} does \emph{worst}, because on big maps the
bottleneck is travel time, not decay.

\paragraph{A corrected false positive.} The table suggested \texttt{mission-aware}
and \texttt{reward-distance} differed on two maps, but inspection of the source
shows their utility reduces to the same function when no mission is active, and
the logs confirm identical behaviour (zero mission intentions). The apparent
gap was small-sample noise --- a reminder not to over-read $n{=}5$
significance.

\paragraph{A diagnosis-driven variant (preliminary).} The diagnosis pinpointed
a scale mismatch: \texttt{reward-distance} scores a pickup by the new parcel
only, while its deliver utility already includes the carried value, biasing it
toward early delivery. A post-hoc \texttt{reward-distance-total} variant scores
a pickup by the whole load's delivered value. In a preliminary smoke test
($n{=}3$) it roughly \emph{doubles} the batch size (26c1\_7: 2.7\,$\to$\,5.3
parcels/delivery; 26c1\_8: 1.3\,$\to$\,2.2) and improves \texttt{reward-distance}
(26c1\_7: 249\,$\to$\,368; 26c1\_8: 409\,$\to$\,537), confirming the diagnosis
--- but it still does not beat \texttt{greedy}. The original strategies and
defaults were left unchanged; this is an experimental variant, not a redesign.

\paragraph{When greedy works, and where it would not.} \texttt{greedy} wins
because Challenge~1 rewards throughput (near-uniform rewards, frequent spawns).
Its expected weaknesses --- heterogeneous rewards, aggressive decay, tight
capacity, contested parcels, one-way layouts, or external missions --- are
largely absent from this dataset; the only non-dominant maps are 26c1\_3 and
26c1\_5 (both ties), and the natural arena for a non-greedy strategy is
Challenge~2, where compliance with missions, not farming, decides the score.

\paragraph{Live compatibility check.} A run against the course server
(\texttt{deliveroojs.bears.disi.unitn.it}) confirmed connection, token
issuance, map loading and logging all work, and that pickup acks carry no ids
there either (validating the reconciliation layer). The live map was
\texttt{crates\_one\_way} (9$\times$9); the agent found most tiles unreachable
and scored~0 --- a compatibility limit (no crate modeling), not a strategy
benchmark.

\paragraph{Challenge~2 (Agent~B) validation.} On the local server with the real
course LLM gateway (\texttt{llama-3.3-70b} over VPN, so \texttt{source=llm} rows
are genuine interpretations) and the official mission agents connecting as god,
Agent~B passed all eight level~1--2 scenarios (Table~\ref{tab:c2}). Seven are
LLM-driven; 26c2\_9 (red/green light) is fallback by design, since the gate must
be instantaneous. The two forbidden-mission scenarios (26c2\_4, 26c2\_6) are
passes in the sense of \emph{no penalty}: the pre-apply (Section~\ref{sec:llm})
blocks the forbidden tiles before the agent can cross them. On 26c2\_4 a pre-fix
run took the $-377$ penalty when an 8.7\,s interpretation arrived too late;
with the pre-apply the same scenario incurs no penalty ($+445$ in a re-run,
$+648$ in the full suite). The PDDL on/off comparison is reported separately in
Table~\ref{tab:pddl} (Section~\ref{sec:pddl}).

\begin{table}[ht]
\centering\small
\begin{tabular}{@{}llll@{}}
\toprule
Scenario & Mission & Source & Outcome \\
\midrule
26c2\_3 & question-answer (calc $\to$ 22)      & llm      & PASS, $+10000$ \\
26c2\_1 & \texttt{go\_to} bonus                & llm      & PASS, $+1000$ \\
26c2\_2 & \texttt{deliver\_at} (1,1)           & llm      & PASS, $+1000$ \\
26c2\_5 & \texttt{deliver\_exactly\_n} (3)     & llm      & PASS, $+100$ \\
26c2\_7 & \texttt{deliver\_less\_value\_than} ($\leq$10) & llm & PASS, $+1000$ \\
26c2\_4 & \texttt{go\_to} \emph{forbidden}     & llm      & PASS, no penalty ($+648$) \\
26c2\_6 & \texttt{deliver\_at} \emph{forbidden}& llm      & PASS, no penalty ($+359$) \\
26c2\_9 & red/green light                      & fallback & PASS, no penalty \\
\bottomrule
\end{tabular}
\caption{Challenge~2 levels~1--2, campaign \texttt{c2-full-v5} (one run per
scenario; the timing-sensitive passes were re-confirmed once). Local server,
official mission agents, real LLM gateway. Honest caveat: encouraging, not a
five-run statistic --- but the failure modes are deterministic in code.}
\label{tab:c2}
\end{table}
